\section{Flat frequency arrays}\label{sec:flat}

Theorem~\ref{thm:frame} yields a linear distribution estimate in one
space dimension.  The useful output measure has a stronger bound along
moving tubes than a general two-dimensional measure.  Projecting a
three-dimensional measure from a slab of width $L$ produces exactly this
bound.  These two reductions account for the exponents $L^{-5/6}$ and
$L^{-4/3}$ below.

\subsection{A one-dimensional distribution estimate}

Fix a bounded closed interval $I\subset\R$ of positive length.  A
one-dimensional packet array at scale $L\geq1$ has the form
\begin{equation}\label{eq:flat-oned-field}
 V(x,t)=L^{-1/2}\sum_{\ell}c_\ell
 e^{i(x\xi_\ell-t\xi_\ell^2)}
 H_\ell\left(\frac{x-b_\ell-2\xi_\ell t}{L},
                  \frac{t}{L^2}\right),\qquad \xi_\ell\in I.
\end{equation}
The array is finite, the pairs $(b_\ell/L,L\xi_\ell)$ have a fixed
bounded number in every unit ball of $\R^2$, and the profiles have fixed
compact support and uniformly bounded derivatives through a sufficiently
large fixed order.  The centers $b_\ell$ are arbitrary real numbers;
there is no integrality assumption.  Constants below may depend on these
fixed bounds, on $I$, and on the fixed size of the output box.

\Needspace{12\baselineskip}
\begin{proposition}[One-dimensional flat estimate]\label{prop:oned}
Let $\nu$ be a nonnegative measure supported in a spatial interval of
length $O(L^2)$ and in $|t|\leq L^2$.  Suppose that, for every
$t_0,b\in\R$, every $v\in I$, and every $1\leq\rho\leq r\leq L^2$,
\begin{equation}\label{eq:flat-tube-test}
 \nu\bigl\{(x,t): |t-t_0|\leq r,
                 \ |x-b-2vt|\leq\rho\bigr\}
       \leq m_1 r\left(1+\frac{\rho}{L}\right).
\end{equation}
There is an exponent $p_0>2$, for which one may take $p_0=21/10$, such
that every array \eqref{eq:flat-oned-field} and every $g>0$ satisfy
\begin{equation}\label{eq:src9}
 \nu\{|V|\geq gL^{-5/6}\}
       \leq C m_1 L^3 g^{-p_0}\|c\|_2^{p_0}.
\end{equation}
The constant is uniform in $L$, in the number of packets, and in their
real centers.  In particular, no factor growing with $L$ is allowed.
\end{proposition}

\begin{proof}
The zero-coefficient case is immediate.  Homogeneity reduces the proof
to $\sum_\ell|c_\ell|^2=1$, but we retain the homogeneous form
\eqref{eq:src9} in recursive applications to subarrays.
We will obtain a broad bound with constant depending on $K$ and a
narrow contraction with a constant independent of $K$.  We then
choose $K$, the base scale, and the induction constant in that order.

First, the hypotheses give two elementary bounds:
\begin{equation}\label{eq:flat-elementary}
 \nu(\R^2)\leq C m_1L^3,\qquad \|V\|_\infty\leq C
                    \quad\text{on } |t|\leq L^2.
\end{equation}
For the mass bound, cover the output box by a fixed number of tests
\eqref{eq:flat-tube-test} with $r=\rho=L^2$ and any fixed $v\in I$.
For the field bound, a packet contributing at $(x,t)$ must satisfy
$|b_\ell/L+2(t/L^2)L\xi_\ell-x/L|\leq C$.  Because $\xi_\ell\in I$
and $|t|/L^2\leq1$, this part of phase space is covered by $O(L)$ unit
balls.  At most $O(L)$ labels contribute, and Cauchy--Schwarz cancels
the normalization $L^{-1/2}$.  It follows that a nonempty level set has
\begin{equation}\label{eq:flat-g-range}
 g\leq C L^{5/6}.
\end{equation}
The mass bound settles $g<1$, and also settles all bounded scales once
the induction constant is chosen sufficiently large.  The same tube
test, with $r=\rho=1$ and a fixed bounded covering when $v\neq0$,
shows that every unit lattice square has $\nu$-mass at most $C m_1$.

\paragraph{The broad alternative.}
Choose an integer $K\geq4$, eventually fixed sufficiently large, and
partition $I$ into $K$ equal half-open intervals $I_j$.  Make a
consistent endpoint assignment at the two ends of $I$.  Let $V_j$ be
the field of the labels in $I_j$, and write $A=gL^{-5/6}$.  Call a
point of $\{|V|\geq A\}$ broad if there are indices $j_1,j_2$ with
$|j_1-j_2|\geq3$ and
\[
 |V_{j_i}(x,t)|\geq \frac{A}{10K},\qquad i=1,2.
\]
The two frequency intervals have length $\delta=|I|/K$ and separation
at least $2\delta$.  We estimate each pair and then sum over at most
$K^2$ pairs; all constants in this paragraph may depend on $K$.

To use the normalization of Theorem~\ref{thm:frame}, put
$W_i=LV_{j_i}$.  Thus the frame amplitudes are
$s_\ell=L^{1/2}c_\ell$, the frame scale is $h=L$, and
\begin{align}
 S_i(x,t)^2
   &=L\sum_{\xi_\ell\in I_{j_i}}|c_\ell|^2
     \left(1+\frac{|x-b_\ell-2\xi_\ell t|}{L}
                  +\frac{|t|}{L^2}\right)^{-P},
       \label{eq:flat-frame-envelopes}\\
 H_i^*&=L^3\sum_{\xi_\ell\in I_{j_i}}|s_\ell|^2\leq L^4,
 \qquad |W_i|\geq A_i:=\frac{gL^{1/6}}{10K}.
       \label{eq:flat-frame-levels}
\end{align}
Here $P$ is a sufficiently large fixed envelope power as in
Theorem~\ref{thm:frame}.  Multiplying the fields by $L$ is essential:
the square envelopes in the following estimates belong to $W_i$.

We claim
\begin{equation}\label{eq:flat-envelope-integrals}
 \int S_i^2\,d\nu\leq C m_1L^3,
 \qquad
 \int S_1^2S_2^2\,d\nu\leq C_Km_1L^3.
\end{equation}
A width-$L$ tube over the given time range has mass $O(m_1L^2)$ by
\eqref{eq:flat-tube-test}.  Two such tubes with velocities separated
by at least $\delta$ intersect only in a time interval of length
$O_K(L)$; their intersection therefore has mass $O_K(m_1L)$.
These statements also control the tails in
\eqref{eq:flat-frame-envelopes}.  More explicitly, a width-$aL$ tube,
with $a\geq1$, has mass at most $C m_1L^2(1+a)$, and the intersection
of width-$aL$ and width-$bL$ tubes at the separated velocities has
mass at most $C_Km_1L(1+a+b)^2$.  To see the latter, subtract the two
tube equations to obtain a time interval of length $O_K((a+b)L)$,
then apply \eqref{eq:flat-tube-test} to the narrower tube.  If this
time length or width exceeds $L^2$, use the original time restriction
and, if necessary, the total mass bound in
\eqref{eq:flat-elementary}; the same polynomial bound results.
Dyadic enlargement and the decay $2^{-Pj}$ absorb these polynomial
costs.  Multiplying the individual tube bound by
$L\sum|c_\ell|^2\leq L$, and the intersection bound by
$L^2(\sum_{\xi_\ell\in I_{j_1}}|c_\ell|^2)
      (\sum_{\xi_\ell\in I_{j_2}}|c_\ell|^2)\leq L^2$,
proves \eqref{eq:flat-envelope-integrals}.

Take $q=9/4$.  For $g\geq1$, discard the points where
$S_i>g^{q/2}=g^{9/8}$ for either $i$.  The first bound in
\eqref{eq:flat-envelope-integrals} charges these points at most
\begin{equation}\label{eq:flat-square-discard}
 C m_1L^3g^{-q}.
\end{equation}
On what remains, use dyadic classes
$b_i\leq S_i<2b_i$, with $b_i\leq g^{9/8}$.
At a witness in such a class the frame upper thresholds can be taken
as $B_i=2b_i$.  By \eqref{eq:flat-g-range},
\begin{equation}\label{eq:flat-frame-ratio}
 \frac{A_i}{B_i}
    \geq c_K L^{1/6}g^{-1/8}
    \geq c'_K L^{1/16}.
\end{equation}
Thus, for $L$ exceeding a fixed bound depending on $K$, both the
large-ratio hypothesis in Theorem~\ref{thm:frame} and its condition
$\delta\geq L^{-1/2}$ hold.  All smaller scales belong to the
induction base.  An envelope cannot vanish at a witness with
$|W_i|>0$, so these positive dyadic classes cover the retained set.

Apply Theorem~\ref{thm:frame} to count the unit squares having a
witness in a fixed class.  Their masses are bounded by $C m_1$, even
when the witness is not a lattice point.  Equations
\eqref{eq:flat-frame-levels}--\eqref{eq:flat-frame-ratio} give the bound
\[
 C_Km_1\frac{b_1^2b_2^2}{A_1^2A_2^2}
       \left(\frac{L^4}{A_1^2}+\frac{L^4}{A_2^2}\right)
   \leq C_Km_1L^3(b_1b_2)^2g^{-6}.
\]
On the other hand, the second bound in
\eqref{eq:flat-envelope-integrals} gives
$C_Km_1L^3(b_1b_2)^{-2}$ for the mass of the class.  Consequently its
mass is at most
\begin{equation}\label{eq:flat-class-minimum}
 C_Km_1L^3\min\bigl\{(b_1b_2)^{-2},
                              (b_1b_2)^2g^{-6}\bigr\}.
\end{equation}
The minimum is bounded by $g^{-3}$.  There are
$O(\log^2(2+g))$ classes with both $b_i\geq g^{-10}$.
For the other classes use the second term in
\eqref{eq:flat-class-minimum} and sum geometrically: for example,
\[
 g^{-6}\sum_{b_1<g^{-10}}b_1^2
                 \sum_{b_2\leq g^{9/8}}b_2^2
       \leq C g^{-6-20+9/4}\leq C g^{-3}.
\]
The analogous sum handles $b_2<g^{-10}$.  Including
\eqref{eq:flat-square-discard} and the choice of the pair proves
\begin{equation}\label{eq:flat-broad-bound}
 \nu(\text{broad points})
   \leq C_Km_1L^3
       \bigl(g^{-9/4}+g^{-3}\log^2(2+g)\bigr)
   \leq C_Km_1L^3g^{-p_0},
 \qquad p_0=21/10.
\end{equation}

\paragraph{The narrow alternative and the exact child hypotheses.}
At a point which is not broad, call $I_j$ significant if
$|V_j|\geq A/(10K)$.  All significant indices are within two steps
of each other.  Choose three consecutive whole partition intervals
containing every significant bin, and let $J$ be the closed hull of
those three intervals.  Let $V_J$ be the sum of exactly their three
bin fields, with the original
endpoint assignments.  The list of all such clusters lies inside
$I$, and each label belongs to at most three clusters.  Near an
endpoint use the first or last three intervals.  Every bin outside
the selected cluster is insignificant; hence
\begin{equation}\label{eq:flat-narrow-level}
 |V_J|\geq |V|-\sum_{I_j\not\subset J}|V_j|
                 \geq\frac9{10}A.
\end{equation}
In particular, the lower fraction in this narrow alternative is
independent of $K$.  Using whole partition intervals avoids dropping
only part of any field at a cluster boundary.

Put $u=3/K$ and choose $v_J$ so that $J=v_J+uI$.
The change of variables and packet parameters are
\begin{equation}\label{eq:flat-child-change}
 x'=u(x-2v_Jt),\quad t'=u^2t,\quad L'=uL,
 \quad b'_\ell=ub_\ell,\quad
 \xi'_\ell=\frac{\xi_\ell-v_J}{u}\in I.
\end{equation}
Up to a common factor of modulus one,
\begin{equation}\label{eq:flat-child-normalization}
 V_J(x,t)=u^{1/2}V'_J(x',t'),
\end{equation}
where $V'_J$ has exactly the form \eqref{eq:flat-oned-field} at scale
$L'$, with unchanged profiles and coefficients.  Indeed,
$b'_\ell/L'=b_\ell/L$ and
$L'\xi'_\ell=L\xi_\ell-Lv_J$, so phase-space ball counts are
unchanged.  This observation also explains why real centers cause
no difficulty in the induction.

Let $\nu'$ be the pushforward of $\nu$ under
\eqref{eq:flat-child-change}.  An arbitrary child test, with parameters
$t'_0,b'\in\R$, $v'\in I$, and
$1\leq\rho'\leq r'\leq(L')^2$, pulls back to
an old test with
\[
 t_0=\frac{t'_0}{u^2},\qquad b=\frac{b'}u,\qquad
 v=v_J+uv'\in J\subset I,\qquad
 r=\frac{r'}{u^2},\qquad \rho=\frac{\rho'}u.
\]
Here $1\leq\rho\leq r\leq L^2$ and
$\rho/L=\rho'/L'$.  Thus the full hypothesis
\eqref{eq:flat-tube-test}, including its lower scale cutoff and its
range of velocities, survives with
\begin{equation}\label{eq:flat-child-measure}
 m'_1=u^{-2}m_1.
\end{equation}
There is no density Jacobian to insert: $\nu'$ is a pushforward
measure.  The time restriction becomes $|t'|\leq(L')^2$.

The transformed spatial domain can be longer than a single interval
of length $(L')^2$.  Partition it into such intervals.  In each one
retain only child packets whose compact support meets that interval
during $|t'|\leq(L')^2$.  The field agrees with this subarray on the
restricted measure.  Because the child velocities are in the fixed
interval $I$ and $L'\geq1$, each compact child packet meets only a
fixed number of these spatial intervals.  If $e_{J,Q}$ is the
coefficient norm of the subarray for cluster $J$ and spatial interval
$Q$, bounded spatial multiplicity and the threefold frequency
overlap give
\begin{equation}\label{eq:flat-child-energy}
 \sum_{J,Q}e_{J,Q}^2\leq C\|c\|_2^2=C,
 \qquad
 \sum_{J,Q}e_{J,Q}^{p_0}
     \leq\left(\sum_{J,Q}e_{J,Q}^2\right)^{p_0/2}\leq C.
\end{equation}
These constants are independent of $K$, even though the number of
spatial intervals can grow with $K$.

By \eqref{eq:flat-narrow-level} and
\eqref{eq:flat-child-normalization}, the child level parameter is
\[
 g'=\frac9{10}g u^{1/3},
 \qquad \frac{(9/10)gL^{-5/6}}{u^{1/2}}
                  =g'(L')^{-5/6}.
\]
Suppose the homogeneous estimate is known at the smaller scale $L'$
with induction constant $C_{\mathrm{ind}}$.  Sum it over the clusters
and spatial intervals, using \eqref{eq:flat-child-measure} and
\eqref{eq:flat-child-energy}.  The narrow contribution is at most
\begin{align}
 C_{\mathrm{ind}}\sum_{J,Q}
       m'_1(L')^3(g')^{-p_0}e_{J,Q}^{p_0}
 &\leq C_0 C_{\mathrm{ind}}
          u^{1-p_0/3}m_1L^3g^{-p_0}.
        \label{eq:flat-narrow-contraction}
\end{align}
The constant $C_0$ does not depend on $K$.  Since
$1-p_0/3=3/10>0$, fix $K$ sufficiently large that
$C_0(3/K)^{3/10}\leq1/2$.  Next fix a base scale $L_*(K)$ large
enough for \eqref{eq:flat-frame-ratio}, the frame separation condition,
and $uL\geq1$ whenever $L>L_*(K)$.  Choose
$C_{\mathrm{ind}}$ to cover the bounded-scale estimate and twice the
constant in \eqref{eq:flat-broad-bound}.  Induction over successive
scale ranges with ratio $u^{-1}$ now closes, proving
\eqref{eq:src9} for all real $L\geq1$.
\end{proof}

\subsection{Projection of a thin frequency strip}

\begin{proposition}[Thin-strip estimate]\label{prop:strip}
Let $U$ be a finite two-dimensional packet array of the form
\eqref{eq:src5}, with the fixed profile and phase-space bounds stated
there.  Suppose all frequency labels lie within distance $C/L$ of an
affine line in frequency space.  Let $\mu$ be supported in a spatial
square of side $O(L^2)$ and in $|t|\leq L^2$, and suppose
$\mu(B_r)\leq mr^2$ for every ball with arbitrary center and every
$r\geq1$.  Then, with the exponent $p_0$ of
Proposition~\ref{prop:oned},
\begin{equation}\label{eq:src10}
 \mu\{|U|\geq gL^{-4/3}\}
       \leq C mL^4g^{-p_0}\|c\|_2^{p_0},\qquad g>0.
\end{equation}
The constant depends only on the fixed bounds, including the strip
width constant.  No sparsity hypothesis is required.
\end{proposition}

\begin{proof}
Rotate the spatial and frequency coordinates so that the frequency
line has constant normal coordinate.  Subtract this central normal
frequency by the corresponding Galilean shear and modulation.
These transformations preserve the field modulus, up to composing
with the change of variables, and change the ball bound by a fixed
factor: the spatial rotation is orthogonal and the space-time shear
and its inverse have bounded norms.  They also preserve uniform
profile and phase-space bounds.  In the new coordinates write the
spatial variables as $(x,y)$ and the frequency labels as
$(\xi_\alpha,\eta_\alpha)$, where
\begin{equation}\label{eq:flat-normal-width}
 |L\eta_\alpha|\leq C.
\end{equation}
The packet fields, including their coefficients, are now
\[
 L^{-1}c_\alpha
 e^{i(x\xi_\alpha+y\eta_\alpha
                     -t(\xi_\alpha^2+\eta_\alpha^2))}
 H_\alpha\left(
       \frac{x-b_{\alpha,x}-2\xi_\alpha t}{L},
       \frac{y-b_{\alpha,y}-2\eta_\alpha t}{L},
       \frac{t}{L^2}\right).
\]
We keep $m$ for the ball parameter after its fixed enlargement.

Partition the normal coordinate into half-open slabs of width $L$,
and let $y_s$ denote the center of slab $s$.  By
\eqref{eq:flat-normal-width} and compact profile support, only labels
in a set $\mathcal A_s$ with
\begin{equation}\label{eq:flat-normal-intercepts}
 |b_{\alpha,y}-y_s|\leq C L
\end{equation}
can contribute in this slab during $|t|\leq L^2$.  Every label
belongs to only boundedly many such sets.  Moreover, the projected
pairs $(b_{\alpha,x}/L,L\xi_\alpha)$ have bounded counts in unit
balls for $\alpha\in\mathcal A_s$.  Indeed, over any such ball, the
remaining coordinates $(b_{\alpha,y}/L,L\eta_\alpha)$ lie in a
bounded box by \eqref{eq:flat-normal-width} and
\eqref{eq:flat-normal-intercepts}.  Its center may depend on $s$,
but it is covered by a fixed number of unit balls.  The original
four-dimensional phase-space bound therefore bounds the number of
projected labels uniformly.  Projection does not cost a power of $L$.

Let $\nu_s$ be the pushforward to $(x,t)$ of the restriction of the
transformed measure to slab $s$.  For any bounded tangential velocity
$v$ and $1\leq\rho\leq r\leq L^2$, the inverse image of a test in
\eqref{eq:flat-tube-test} has dimensions comparable to $r$, $\rho$,
and $L$, after a shear of bounded norm.  It can be covered by
\[
 C\frac r\rho\left(1+\frac L\rho\right)
\]
balls of radius $C\rho$.  These radii are at least one.  Applying the
ball bound gives
\begin{equation}\label{eq:flat-projected-test}
 \nu_s\{|t-t_0|\leq r,\ |x-b-2vt|\leq\rho\}
      \leq Cmr(L+\rho)
      =CmL\,r\left(1+\frac\rho L\right).
\end{equation}
Thus Proposition~\ref{prop:oned} applies to $\nu_s$ with
$m_1=CmL$ and any fixed tangential frequency interval containing
the labels.

It remains to handle the dependence of the field on the normal
coordinate.  A pointwise choice of a normal slice would not suffice
for an arbitrary measure, so we use a single Fourier expansion
valid throughout the slab.  Put
\[
 z=\frac{x-b_{\alpha,x}-2\xi_\alpha t}{L},\qquad
 \tau=\frac{t}{L^2},\qquad w=\frac{y-y_s}{L},\qquad
 \zeta_{\alpha,s}=\frac{y_s-b_{\alpha,y}}L,
 \quad a_\alpha=L\eta_\alpha.
\]
Choose a fixed smooth cutoff $\chi(w)$ equal to one for
$|w|\leq1/2$ and supported strictly inside one period of fixed length
$T$.  Expand
\begin{align}
 F_{\alpha,s}(z,\tau,w)
   &:=\chi(w)e^{i a_\alpha w-i a_\alpha^2\tau}
       H_\alpha(z,w+\zeta_{\alpha,s}-2a_\alpha\tau,\tau)
       \notag\\
   &=\sum_{n\in\mathbb Z}e^{2\pi i n w/T}
                          F_{\alpha,s,n}(z,\tau).
       \label{eq:flat-normal-series}
\end{align}
The constants $a_\alpha$ and $\zeta_{\alpha,s}$ are bounded on
$\mathcal A_s$.  Consequently, if $D$ is the finite derivative order
needed in Proposition~\ref{prop:oned}, integration by parts in $w$
gives, for any sufficiently large fixed $N$ for which the original
profiles have $D+N$ derivatives,
\begin{equation}\label{eq:flat-mode-decay}
 \max_{j+k\leq D}
  \|\partial_z^j\partial_\tau^kF_{\alpha,s,n}\|_\infty
       \leq C_N(1+|n|)^{-N}.
\end{equation}
All these profiles have support in a common bounded $(z,\tau)$ box.
The residual normal phase has bounded derivatives in these scaled
variables; this is the reason for first removing the central normal
frequency.  The remaining constant phase
$e^{iy_s\eta_\alpha}$ can be included in the coefficient.

Absorb $C_N(1+|n|)^{-N}$ from \eqref{eq:flat-mode-decay} into the
coefficients and normalize each remaining profile.  This produces
one-dimensional fields $V_{s,n}$ with uniform admissible profiles and
the already verified projected phase-space counts, such that in
slab $s$
\begin{equation}\label{eq:flat-strip-series}
 U(x,y,t)=L^{-1/2}\sum_{n\in\mathbb Z}
       e^{2\pi i n (y-y_s)/(TL)}V_{s,n}(x,t),
 \qquad
 \|c_{s,n}\|_2\leq C_N(1+|n|)^{-N}e_s,
\end{equation}
where $e_s^2=\sum_{\alpha\in\mathcal A_s}|c_\alpha|^2$.
The series converges absolutely, and all its external phase factors
have modulus one at every witness.  The factor $L^{-1/2}$ outside
the series is the difference between the two-dimensional packet
normalization $L^{-1}$ and the one-dimensional normalization
$L^{-1/2}$.

Choose positive numbers $a_n=c_*(1+|n|)^{-2}$ with
$\sum_n a_n=1/2$.  If $|U(x,y,t)|\geq gL^{-4/3}$ in slab $s$,
then \eqref{eq:flat-strip-series} forces, for at least one $n$,
\[
 |V_{s,n}(x,t)|\geq a_ngL^{-5/6}.
\]
Taking the projected measure of this union and applying
Proposition~\ref{prop:oned} gives
\begin{align}
 \mu\bigl(\{|U|\geq gL^{-4/3}\}\cap\text{slab }s\bigr)
 &\leq\sum_{n\in\mathbb Z}
       \nu_s\{|V_{s,n}|\geq a_ngL^{-5/6}\}
       \notag\\
 &\leq C mL^4g^{-p_0}e_s^{p_0}
       \sum_{n\in\mathbb Z}a_n^{-p_0}(1+|n|)^{-Np_0}
       \notag\\
 &\leq C mL^4g^{-p_0}e_s^{p_0},
       \label{eq:flat-mode-sum}
\end{align}
provided $(N-2)p_0>1$.  Only finitely many derivative bounds are
needed to choose such an $N$.  Finally,
\[
 \sum_s e_s^2\leq C\|c\|_2^2,
 \qquad
 \sum_s e_s^{p_0}
       \leq\left(\sum_s e_s^2\right)^{p_0/2}
       \leq C\|c\|_2^{p_0},
\]
by bounded slab overlap and $p_0>2$.  Summing
\eqref{eq:flat-mode-sum} proves \eqref{eq:src10}.
\end{proof}

\begin{remark}[Fixed enlargements of the time domain]\label{rem:flat-time}
The preceding conclusions also hold on $|t|\leq C_tL^2$ for any fixed
$C_t$, with constants depending on $C_t$.  Partition this range into
boundedly many intervals of radius at most $L^2$, and translate the
center of each interval to zero.  A time translation by $t_*$ replaces
the packet intercept $b$ by $b+2\omega t_*$ and shifts the profile time
argument by $t_*/L^2$; fixed support and derivative bounds remain
uniform.  On phase space this is the bounded shear
$(b/L,L\omega)\mapsto(b/L+2(t_*/L^2)L\omega,L\omega)$,
so real-center ball counts remain bounded.  The one-dimensional
tube hypothesis is invariant under this translation because its
time centers and intercepts are arbitrary; the two-dimensional ball
hypothesis is translation invariant as well.  Restriction to each
time interval and a bounded sum therefore give the asserted extension.
\end{remark}
